\documentclass[14pt,a4paper]{extarticle}
\usepackage[T2A]{fontenc}
\usepackage[utf8]{inputenc}
\usepackage[english,russian]{babel}
\usepackage{lmodern}
\usepackage{geometry}
\usepackage{array}
\usepackage{tabularx}
\usepackage{longtable}
\usepackage{ltablex}
\usepackage{booktabs}
\usepackage{hyperref}
\usepackage{indentfirst}
\usepackage{microtype}
\usepackage{graphicx}
\usepackage{setspace}
\usepackage{listings}
\usepackage{graphicx}
\usepackage{caption}
\usepackage{float}
\usepackage{tocloft}
\usepackage{titlesec}
\usepackage{newtxtext,newtxmath}

\titleformat{\section}        % форматируем \section
  {\normalfont\bfseries}      % стиль: обычный + жирный
  {\thesection.}              % номер секции с точкой
  {2mm}                       % расстояние между номером и текстом
  {}                           % перед заголовком можно добавить что-то
\titlespacing*{\section}{1.5cm}{*3}{*1}  % {отступ слева}{отступ сверху}{отступ снизу}

\captionsetup{justification=centering}
\renewcommand{\ttdefault}{cmtt}
\keepXColumns
\setlength{\parindent}{3em}
\newgeometry{left=2.5cm,right=2.5cm,top=3cm,bottom=2.5cm}

\renewcommand{\contentsname}{\hspace*{4em}{\smallСодержание}}
\setlength{\cftbeforesecskip}{0.5em}
\renewcommand{\cftsecleader}{\cftdotfill{\cftdotsep}}
\renewcommand{\cftsecaftersnum}{.\hspace{2mm}}

\title{Отчёт по "Организация работы ИТ-службы"}
\author{Выполнил студент группы М-РИТ-25-1, Станиславчук Сергей}

\begin{document}

\setstretch{1.5}
\noindent
\begin{minipage}[t]{0.01\textwidth}
    \includegraphics{~/programmer/++Programmer/5_1/LSTU.jpg}
\end{minipage}%
\begin{minipage}[c]{1.1\linewidth}
    \centering
    {\small
    МИНИСТЕРСТВО НАУКИ И ВЫСШЕГО ОБРАЗОВАНИЯ РФ \\[1em]
    ФЕДЕРАЛЬНОЕ ГОСУДАРСТВЕННОЕ БЮДЖЕТНОЕ \\
    ОБРАЗОВАТЕЛЬНОЕ УЧРЕЖДЕНИЕ ВЫСШЕГО ОБРАЗОВАНИЯ \\
    «ЛИПЕЦКИЙ ГОСУДАРСТВЕННЫЙ ТЕХНИЧЕСКИЙ УНИВЕРСИТЕТ»
    }
\end{minipage}%

\vspace{3em}

\noindent
\hspace{1em}%
\begin{minipage}[t]{0.8\textwidth}
\textbf{Институт} \hspace{2em} компьютерных наук \\
\textbf{Кафедра} \hspace{2em} автоматизированных систем управления
\end{minipage}


\vspace{6em}

\begin{center}
\textbf{Практическая работа №4} \\
по дисциплине \\ <<Организация работы ИТ-службы>>
\end{center}

\vspace{6em}

\noindent
\begin{tabular*}{\textwidth}{@{\extracolsep{\fill}} l c r }
    Студент М-РИТ-25-1 & \makebox[4cm]{\hrulefill} & Станиславчук С.М. \\[2em]
Руководитель & \\
К.т.н., доцент & \makebox[4cm]{\hrulefill} & Ведищев В.В. \\
\end{tabular*}

\vfill
\begin{center}
    Липецк, 2026
\end{center}

\thispagestyle{empty}

\newpage
\tableofcontents
\newpage

\section{Координация проектирования}

Координация проектирования обеспечивает согласованность работ, выполняемых при
создании или изменении IT-сервиса. Для VDS-хостинга данный процесс необходим,
поскольку предоставление виртуальных серверов зависит от совместной работы
физической инфраструктуры, платформы виртуализации, сетевых компонентов,
систем хранения данных, средств мониторинга и персонала, отвечающего за
эксплуатацию.

Основная задача координации заключается в том, чтобы все изменения сервиса
соответствовали установленным требованиям заказчиков, согласованным параметрам
качества и возможностям поставщика. Например, при запуске нового тарифа с
увеличенным объемом оперативной памяти необходимо проверить наличие свободных
ресурсов на физических узлах, достаточность сетевой и дисковой производительности,
возможность мониторинга новых конфигураций и соответствие изменений действующим
процедурам информационной безопасности.

В рамках координации проектирования выполняются следующие действия:
\begin{itemize}
    \item определение целей и границ проектируемых изменений;
    \item согласование требований между подразделениями эксплуатации,
    технической поддержки, информационной безопасности и управления услугой;
    \item оценка необходимых технических, финансовых и кадровых ресурсов;
    \item определение рисков, зависимостей и возможного влияния изменений на
    действующих заказчиков;
    \item согласование сроков реализации, плана тестирования и критериев приемки;
    \item передача результатов проектирования в эксплуатацию вместе с
    необходимой технической документацией.
\end{itemize}

Например, при модернизации платформы виртуализации необходимо заранее определить,
можно ли выполнить обновление без остановки виртуальных серверов. Если это
невозможно, должны быть согласованы окно технического обслуживания, порядок
уведомления заказчиков и план возврата к предыдущей версии в случае возникновения
критических ошибок.

Результатом координации проектирования является согласованный план изменений,
в котором определены ответственные лица, необходимые ресурсы, сроки, риски,
требования к качеству и порядок передачи результата в эксплуатацию.

\section{Управление каталогом услуг}

Каталог услуг представляет собой структурированное описание услуг, которые
поставщик предоставляет заказчикам. Для VDS-хостинга каталог должен содержать
актуальную информацию о доступных конфигурациях виртуальных серверов, составе
услуги, условиях предоставления и дополнительных возможностях.

Основной услугой является предоставление виртуального сервера с определенным
количеством виртуальных процессорных ядер, объемом оперативной памяти, дисковым
пространством и сетевым подключением. В зависимости от тарифной модели каталог
может включать несколько конфигураций, отличающихся объемом ресурсов и стоимостью.

Для каждой позиции каталога целесообразно указывать:
\begin{itemize}
    \item наименование и назначение тарифа;
    \item количество виртуальных процессорных ядер;
    \item объем оперативной памяти;
    \item тип и объем дискового пространства;
    \item параметры сетевого подключения и ограничения трафика, если они
    предусмотрены;
    \item доступные операционные системы и способы установки;
    \item стоимость и порядок оплаты;
    \item условия изменения конфигурации;
    \item условия резервного копирования и восстановления;
    \item параметры технической поддержки и доступности;
    \item ограничения допустимого использования сервиса.
\end{itemize}

Помимо основной услуги, каталог может включать дополнительные услуги: резервное
копирование виртуальных серверов, предоставление дополнительных IP-адресов,
увеличение дискового пространства, администрирование операционной системы,
защиту от сетевых атак и помощь при переносе данных.

Каталог должен отражать именно те услуги и параметры, которые доступны заказчику.
Внутренние технические компоненты, например физические серверы, системы хранения
и сетевое оборудование, сами по себе не являются отдельными позициями клиентского
каталога, если поставщик не предоставляет их заказчику как самостоятельную услугу.

Ответственность за актуальность каталога возлагается на владельца услуги или
назначенное подразделение управления услугами. При изменении тарифов,
характеристик инфраструктуры, условий обслуживания или состава дополнительных
услуг информация в каталоге должна своевременно обновляться.

Таким образом, каталог услуг обеспечивает единое представление о предложении
поставщика и позволяет заказчику определить, какая конфигурация VDS соответствует
его потребностям.

\section{Управление уровнем услуг}

Управление уровнем услуг направлено на согласование, контроль и улучшение качества
предоставляемого сервиса. Основным инструментом данного процесса является
соглашение об уровне обслуживания (SLA), в котором фиксируются измеримые
параметры услуги и обязательства поставщика.

Для VDS-хостинга существенное значение имеют доступность виртуального сервера,
производительность инфраструктуры, время реакции технической поддержки и сроки
восстановления после инцидентов. Эти параметры должны быть определены таким
образом, чтобы заказчик мог оценить качество услуги, а поставщик -- проверить
выполнение принятых обязательств.

В качестве целевых показателей могут использоваться:
\begin{itemize}
    \item доступность услуги -- например, 99,99\% за год;
    \item время первичной реакции технической поддержки -- в зависимости от
    приоритета обращения;
    \item целевое время восстановления работоспособности после критического
    инцидента;
    \item допустимые параметры производительности дисковой и сетевой подсистем;
    \item периодичность предоставления отчетов о качестве услуги.
\end{itemize}

Для оценки доступности может применяться следующая формула:

\[
    A = \frac{T_{\text{согласованное}} -
    T_{\text{простоя}}}{T_{\text{согласованное}}} \cdot 100\%,
\]

где $A$ -- доступность услуги, $T_{\text{согласованное}}$ -- согласованное время
предоставления услуги, $T_{\text{простоя}}$ -- время простоя, учитываемое согласно
условиям SLA.

Поставщик должен регулярно собирать данные мониторинга, анализировать обращения
заказчиков и сопоставлять фактические показатели с согласованными значениями.
При выявлении отклонений необходимо определить причины их возникновения и
разработать корректирующие мероприятия.

Например, если доступность VDS систематически оказывается ниже целевого значения,
необходимо проверить частоту отказов физических серверов, состояние сетевой
инфраструктуры, работу систем хранения и эффективность процедур восстановления.

Результатом управления уровнем услуг является поддержание согласованного качества
сервиса, прозрачность обязательств поставщика и наличие объективных данных для
оценки предоставляемой услуги.

\section{Управление доступностью}

Управление доступностью обеспечивает способность IT-сервиса выполнять свои
функции в согласованные периоды времени. Для VDS-хостинга данный процесс имеет
особое значение, поскольку виртуальные серверы могут использоваться для
размещения веб-сайтов, баз данных, корпоративных приложений и других систем,
которые должны работать круглосуточно.

Целевым режимом предоставления услуги является 24/7. Доступность VDS зависит не
только от состояния виртуальной машины, но и от работоспособности всех компонентов,
необходимых для ее функционирования: физических серверов, систем хранения,
сетевого оборудования, платформы виртуализации и электропитания дата-центра.

Для обеспечения доступности поставщик должен применять следующие меры:
\begin{itemize}
    \item круглосуточный мониторинг состояния физических серверов, виртуальных
    машин, сетевых соединений и систем хранения;
    \item резервирование критически важных компонентов инфраструктуры;
    \item использование источников бесперебойного питания и резервных систем
    электроснабжения на площадке размещения оборудования;
    \item своевременную замену неисправных компонентов;
    \item наличие процедур миграции виртуальных машин на исправные физические узлы,
    если это поддерживается архитектурой платформы;
    \item регистрацию и анализ инцидентов, приводящих к недоступности услуги;
    \item планирование технического обслуживания с учетом минимизации влияния на
    заказчиков.
\end{itemize}

Для контроля доступности используются показатели времени простоя, количества
инцидентов, средней продолжительности восстановления и доли времени, в течение
которого услуга была доступна. Мониторинг должен обнаруживать не только полный
отказ виртуального сервера, но и ситуации, когда услуга формально работает, однако
ее производительность существенно ухудшилась.

Необходимо также определить порядок проведения плановых работ. Заказчики должны
получать уведомления о работах, способных повлиять на доступность VDS. При этом
плановые перерывы и исключения из расчета SLA должны быть четко определены в
соглашении об уровне обслуживания.

Таким образом, управление доступностью направлено на сокращение количества и
продолжительности отказов, а также на обеспечение соответствия фактической
доступности согласованным требованиям.

\section{Управление мощностями}

Управление мощностями обеспечивает наличие достаточных ресурсов для предоставления
услуг с требуемым уровнем производительности. В случае VDS-хостинга данный процесс
заключается в планировании и контроле вычислительных, дисковых и сетевых ресурсов
физической инфраструктуры.

К основным ресурсам относятся:
\begin{itemize}
    \item процессорная мощность физических серверов;
    \item оперативная память;
    \item емкость и производительность систем хранения данных;
    \item пропускная способность сетевых каналов;
    \item количество доступных физических узлов для размещения виртуальных машин.
\end{itemize}

Поставщик должен контролировать как общую загрузку инфраструктуры, так и
распределение ресурсов между отдельными физическими серверами. Высокая средняя
загрузка может приводить к снижению производительности виртуальных машин,
увеличению времени отклика и невозможности размещения новых заказов.

Управление мощностями включает несколько направлений. Управление мощностями
бизнеса предполагает прогнозирование спроса на VDS и определение необходимого
масштаба инфраструктуры. Управление мощностями услуг направлено на оценку ресурсов,
которые требуются для выполнения требований SLA. Управление мощностями компонентов
заключается в контроле загрузки процессоров, памяти, накопителей и сетевого
оборудования.

Для прогнозирования используются исторические данные о потреблении ресурсов,
количестве активных виртуальных серверов, динамике подключений новых заказчиков и
сезонных изменениях нагрузки. На основании этих данных поставщик может планировать
закупку оборудования, расширение систем хранения и увеличение пропускной
способности каналов связи.

В качестве контрольных показателей могут использоваться:
\begin{itemize}
    \item средняя и пиковая загрузка процессоров;
    \item доля используемой оперативной памяти;
    \item свободная емкость хранилищ;
    \item показатели производительности дисковой подсистемы;
    \item средняя и пиковая загрузка сетевых каналов;
    \item количество свободных ресурсов для создания новых виртуальных серверов.
\end{itemize}

Для предотвращения нехватки ресурсов могут устанавливаться пороговые значения.
Например, при приближении загрузки хранилища к установленному пределу система
мониторинга формирует предупреждение, после чего ответственные сотрудники
оценивают необходимость расширения емкости.

При этом важно учитывать возможность избыточного выделения ресурсов виртуальным
машинам и конкуренцию за физические ресурсы между несколькими заказчиками.
Планирование должно учитывать особенности технологии виртуализации и допустимый
уровень переподписки ресурсов.

Результатом управления мощностями является поддержание необходимого запаса
ресурсов, предотвращение перегрузок и возможность развития сервиса без
существенного ухудшения качества обслуживания.

\section{Управление непрерывностью IT-услуг}

Управление непрерывностью IT-услуг направлено на обеспечение восстановления
сервиса после серьезных аварий и других событий, способных нарушить его
предоставление на длительное время. Для VDS-хостинга такими событиями могут быть
масштабный отказ оборудования, повреждение систем хранения, длительное отключение
электроэнергии, потеря сетевого подключения или недоступность площадки дата-центра.

В отличие от управления доступностью, которое в основном направлено на
предотвращение и сокращение обычных отказов, управление непрерывностью
рассматривает сценарии серьезных нарушений и заранее определяет порядок
восстановления услуги.

Основой данного процесса является анализ влияния на бизнес (Business Impact
Analysis, BIA). Он позволяет определить последствия длительной недоступности VDS
для заказчиков, выявить критичные компоненты сервиса и установить приоритеты
восстановления.

Для планирования непрерывности могут использоваться два показателя:
\begin{itemize}
    \item RTO (Recovery Time Objective) -- целевое время восстановления услуги
    после нарушения;
    \item RPO (Recovery Point Objective) -- максимально допустимый период потери
    данных, выраженный во времени.
\end{itemize}

Конкретные значения RTO и RPO должны определяться с учетом архитектуры сервиса,
стоимости резервирования и требований заказчиков. Для разных тарифов или
дополнительных услуг эти значения могут отличаться.

Для обеспечения непрерывности VDS-сервиса могут применяться следующие меры:
\begin{itemize}
    \item резервное копирование данных виртуальных серверов;
    \item хранение резервных копий отдельно от основной инфраструктуры;
    \item резервирование сетевых каналов и критически важных компонентов;
    \item подготовка резервных вычислительных мощностей;
    \item документирование процедур восстановления виртуальных машин;
    \item наличие плана действий при отказе дата-центра или крупной части
    инфраструктуры;
    \item регулярное тестирование восстановления из резервных копий.
\end{itemize}

Особое внимание должно уделяться проверке резервных копий. Сам факт успешного
создания копии не гарантирует возможность восстановления данных. Поэтому
необходимо периодически выполнять тестовое восстановление виртуальных машин и
проверять целостность полученных данных.

План непрерывности должен содержать перечень ответственных сотрудников, порядок
эскалации инцидентов, контактные данные внешних поставщиков, последовательность
восстановления компонентов и порядок информирования заказчиков.

Результатом управления непрерывностью является подготовленность поставщика к
крупным авариям и наличие проверенных процедур восстановления VDS-сервиса в
пределах установленных целевых показателей.

\section{Управление информационной безопасностью}

Управление информационной безопасностью направлено на обеспечение
конфиденциальности, целостности и доступности информации, обрабатываемой и
хранимой в рамках IT-сервиса. Для VDS-хостинга безопасность имеет особое значение,
поскольку на одной физической инфраструктуре могут размещаться виртуальные
серверы различных заказчиков.

Одной из основных задач является обеспечение изоляции виртуальных машин.
Поставщик должен предотвращать несанкционированный доступ одного заказчика к
ресурсам, данным и сетевому трафику другого заказчика. Для этого применяются
механизмы изоляции платформы виртуализации, разграничение сетевого доступа и
контроль привилегий.

К основным направлениям управления информационной безопасностью относятся:
\begin{itemize}
    \item управление учетными записями и правами доступа сотрудников к
    инфраструктуре;
    \item использование надежной аутентификации и, при необходимости,
    многофакторной аутентификации;
    \item разграничение доступа к панели управления, гипервизорам, системам
    хранения и резервным копиям;
    \item своевременное обновление платформы виртуализации и системного
    программного обеспечения;
    \item защита сетевой инфраструктуры с помощью межсетевых экранов и средств
    обнаружения атак;
    \item регистрация событий безопасности и анализ подозрительной активности;
    \item защита резервных копий и контроль доступа к ним;
    \item реагирование на инциденты информационной безопасности.
\end{itemize}

Ответственность за безопасность распределяется между поставщиком и заказчиком.
Поставщик отвечает за безопасность физической инфраструктуры, платформы
виртуализации, систем управления и тех компонентов, которые находятся в его зоне
контроля. Заказчик отвечает за настройку операционной системы, установку
обновлений, управление учетными записями внутри своего VDS и защиту размещенных
на нем приложений и данных, если иное не предусмотрено условиями услуги.

Важной частью процесса является управление уязвимостями. Поставщик должен
отслеживать появление критических уязвимостей в используемом программном
обеспечении, оценивать их влияние и планировать установку исправлений. Обновления,
которые могут повлиять на работу виртуальных серверов, должны выполняться в
соответствии с процедурами управления изменениями.

Также должны быть определены правила допустимого использования услуги. Запрещается
применение VDS для проведения атак на сторонние системы, распространения
вредоносного программного обеспечения, фишинга и другой противоправной
деятельности. При выявлении нарушений поставщик может ограничить или приостановить
предоставление услуги в соответствии с договором и правилами обслуживания.

Результатом управления информационной безопасностью является снижение риска
утечки, повреждения или потери данных, предотвращение несанкционированного доступа
и поддержание доверия заказчиков к VDS-сервису.

\section{Управление подрядчиками}

Управление подрядчиками обеспечивает контроль качества услуг, которые поставщик
VDS использует для предоставления собственного сервиса. Поскольку работа
виртуальных серверов зависит от физической инфраструктуры и внешних ресурсов,
поставщик может привлекать сторонние организации для размещения оборудования,
обеспечения электропитания, сетевого подключения и других вспомогательных функций.

К основным подрядчикам и внешним поставщикам могут относиться:
\begin{itemize}
    \item оператор дата-центра, предоставляющий помещения, электропитание,
    охлаждение и физическую охрану;
    \item оператор связи, обеспечивающий подключение к сети Интернет и
    транспортировку сетевого трафика;
    \item поставщики серверного оборудования, систем хранения и сетевых
    компонентов;
    \item разработчики и поставщики программного обеспечения виртуализации,
    систем управления и мониторинга;
\end{itemize}

Управление подрядчиками включает выбор поставщиков, согласование условий
сотрудничества, контроль выполнения договорных обязательств и оценку качества
предоставляемых ресурсов.

При выборе подрядчика необходимо учитывать его надежность, технические
возможности, стоимость услуг, репутацию, географическое расположение объектов,
уровень информационной безопасности и способность выполнять согласованные
требования. Например, при выборе дата-центра важно оценить наличие резервного
электропитания, систем охлаждения, физической охраны и процедур аварийного
реагирования.

В договорах с подрядчиками должны быть зафиксированы:
\begin{itemize}
    \item перечень предоставляемых услуг и границы ответственности сторон;
    \item согласованные показатели качества и доступности;
    \item порядок уведомления о плановых работах и авариях;
    \item сроки реакции и устранения неисправностей;
    \item требования к информационной безопасности и конфиденциальности;
    \item порядок обработки инцидентов и взаимодействия с технической поддержкой;
    \item условия изменения и прекращения договора.
\end{itemize}

Поставщик VDS должен контролировать не только выполнение отдельных договоров, но
и влияние работы подрядчиков на собственный SLA. Например, если внешний оператор
связи обеспечивает сетевое подключение, его показатели доступности должны
учитываться при планировании доступности конечной услуги. При этом обязательства
перед заказчиком остаются ответственностью поставщика VDS в пределах заключенного
с ним соглашения.

Для контроля подрядчиков могут анализироваться количество и продолжительность
отказов, соблюдение сроков устранения неисправностей, выполнение требований
безопасности и соответствие фактического качества договорным условиям. При
систематическом нарушении обязательств необходимо инициировать корректирующие
мероприятия или рассмотреть замену подрядчика.

Таким образом, управление подрядчиками позволяет контролировать внешние
зависимости VDS-сервиса, снижать риски нарушения предоставления услуги и
поддерживать требуемый уровень качества.

\end{document}
